% Generated draft: requires mathematical and source review.
\documentclass[11pt]{article}
\usepackage[margin=25mm]{geometry}
\usepackage{fontspec}
\setmainfont{Latin Modern Roman}
\usepackage{amsmath,amssymb,mathtools,unicode-math}
\setmathfont{Latin Modern Math}
\usepackage{enumitem,needspace,xurl,hyperref}
\setlength{\parindent}{0pt}
\setlength{\parskip}{5pt}
\setlength{\emergencystretch}{3em}
\newcommand{\N}{\mathbb{N}}
\newcommand{\Z}{\mathbb{Z}}
\newcommand{\Q}{\mathbb{Q}}
\newcommand{\R}{\mathbb{R}}
\newcommand{\C}{\mathbb{C}}
\newcommand{\F}{\mathbb{F}}
\newcommand{\A}{\mathbb{A}}
\newcommand{\PP}{\mathbb{P}}
\newcommand{\E}{\mathbb{E}}
\newcommand{\eps}{\varepsilon}
\newcommand{\dd}{\,\mathrm{d}}
\begin{document}
\section*{UnsolvedMath ID: 800851}
\textbf{Source problem code:} catalogue entry
\subsection{Definitions and mathematical statement}
For an integer \(d\geq 2\), an orthonormal basis of \(\C^d\) is a family \(B=(b_1,\ldots,b_d)\) satisfying \(\langle b_j,b_k\rangle=\delta_{jk}\) for the standard Hermitian inner product. Two orthonormal bases \(B=(b_j)_{j=1}^d\) and \(C=(c_k)_{k=1}^d\) are mutually unbiased if
\[
\lvert\langle b_j,c_k\rangle\rvert^2=\frac1d
\qquad\text{for all }j,k\in\{1,\ldots,d\}.
\]
A family of orthonormal bases is simultaneously mutually unbiased if every two distinct bases in the family are mutually unbiased. Let \(\mathrm{MUB}(d)\) be the maximum cardinality of such a family.

The supplied source attributes the bound \(\mathrm{MUB}(d)\leq d+1\), with equality for prime-power dimensions, to the reference labeled \texttt{Bandyopadhyay2002-MUBs}. It describes determining \(\mathrm{MUB}(d)\) in non-prime-power dimensions as open and identifies dimension six as the smallest open instance. These are supplied source assessments, not independently verified literature conclusions.

\textbf{Conjecture 1 (Bandeira et al., as attributed in the supplied record).} One has
\[
\mathrm{MUB}(6)<7.
\]
Equivalently, there do not exist seven orthonormal bases \(B_r=(b_{r,1},\ldots,b_{r,6})\), indexed by \(r\in\{1,\ldots,7\}\), of \(\C^6\) such that
\[
\lvert\langle b_{r,j},b_{s,k}\rangle\rvert^2=\frac16
\quad\text{for all }1\leq r<s\leq7\text{ and }1\leq j,k\leq6.
\]
The conjecture excludes a complete family of seven mutually unbiased bases; it does not specify the exact value of \(\mathrm{MUB}(6)\).
\subsection{Short English statement}
The conjecture states that seven pairwise mutually unbiased orthonormal bases cannot exist in \(\C^6\): vectors from different bases would have to have squared inner-product magnitude \(1/6\). The supplied registry labels this conjecture open; that status has not been independently verified.
\subsection{Sources}
\begin{itemize}
\item Local UnsolvedMath registry: permanent ID \(800851\), source problem code \texttt{catalogue entry}, title ``Randomstrasse101 Open Problems (2025) --- Conjecture 1.'' The supplied record attributes the problem to Bandeira et al., gives proposed year 2026, and labels its status open.
\item Supplied source list: Bandeira et al., \emph{Randomstrasse101: Open Problems of 2025} (2026), Conjecture 1. Supplied URL: \url{https://arxiv.org/abs/2603.29571}. The record reports access on 2026-10-16 and extraction by \texttt{arxiv\_latex}; these are supplied provenance details, not a claim of independent access or verification.
\item Reference label \texttt{Bandyopadhyay2002-MUBs}, cited in the supplied statement for the upper bound and prime-power equality. Full bibliographic details and a reference-specific URL were not supplied.
\item Reference label \texttt{Afonso10L42P}, Open Problem 6.2, cited in the supplied statement regarding the dimension-six problem. Full bibliographic details and a reference-specific URL were not supplied.
\end{itemize}
\end{document}
